\subsection{Springer 对应}\label{ss:Spr corr} 设 $W$ 为 $G$ 的 Weyl 群。1976 年，Springer \cite{Springer} 作出了根本性的观察：$\cohog{*}{\sB_{e}}$ 上存在自然的 $W$ 作用，尽管 $W$ 并不作为簇的自同构作用在 $\sB_{e}$ 上。

\begin{theorem}[Springer {\cite[Thm 6.10]{Springer}}]\label{th:Spr corr}
\begin{enumerate}
\item 对每个幂零元 $e$，$\cohog{*}{\sB_{e}}$ 上有自然的 $W$ 分次作用，且与 $A_{e}$ 的作用交换。
\item 对每个幂零元 $e$ 与 $A_{e}$ 的每个不可约表示 $\rho$，重数空间 $M(e,\rho):=\Hom_{A_{e}}(\rho, \cohog{2d_{e}}{\sB_{e}})$ 或为零，或在 (1) 中的作用之下是 $W$ 的不可约表示。
\item $W$ 的每个不可约表示 $\chi$ 都在 $G$ 共轭意义下唯一地以 $M(e,\rho)$ 的形式出现。于是对应 $\chi\mapsto (e,\rho)$ 给出单射
\begin{equation}\label{Spr corr}
\Irr(W)\incl \{(e,\rho)\}/G.
\end{equation}
\end{enumerate}
\end{theorem}

\subsubsection{约定}\label{sss:two actions} 事实上 $\cohog{*}{\sB_{e}}$ 上有两个自然的 $W$ 作用，二者相差与 $W$ 的符号表示作张量。本讲义采用由以下性质规范的那个作用：$W$ 的平凡表示对应正则幂零元 $e$ 与平凡 $\rho$；$W$ 的符号表示对应 $e=0$。但要注意，Springer 的原始论文 \cite{Springer} 用的是另一个作用。

\subsubsection{$e=0$ 的情形} 取 $e=0$，Springer 定理给出 $W$ 在 $\cohog{*}{\calB}$ 上的分次作用。这是什么作用？首先，可以改用 $G/T$ 而非 $\sB=G/B$ 来几何地看出这个作用。事实上，由于 $N_{G}(T)/T=W$，$N_{G}(T)$ 在 $G/T$ 上的右作用诱导出 $W$ 在 $G/T$ 上的作用，进而诱导出 $W$ 在 $\cohog{*}{G/T}$ 上的作用。由于投影 $G/T\to G/B$ 是仿射空间丛，有 $\cohog{*}{\sB}\cong\cohog{*}{G/T}$。可以证明，在此同构之下，$W$ 在 $\cohog{*}{G/T}$ 上的作用恰对应 Springer 在 $\cohog{*}{\calB}$ 上的作用。

设 $S=\Sym(\xch(T)\otimes\Ql)$ 为分次对称代数，其中 $\xch(T)$ 的次数为 $2$。$W$ 在 $\xch(T)$ 上的反射表示随即诱导出 $W$ 在 $S$ 上的分次作用。回忆旗簇上同调环的 Borel 表示
\begin{equation}\label{Borel HB}
\cohog{*}{\calB,\Ql}\cong S/(S^{W}_{+})
\end{equation}
其中 $S_{+}\subset S$ 是由正次数元素张成的理想，$(S^{W}_{+})$ 表示 $S$ 中由 $S_{+}$ 的 $W$ 不变量生成的理想。其实 \eqref{Borel HB} 是 $W$ 模的同构（见 \cite[Prop 7.2]{Springer}）。由 Chevalley 的定理，$S/(S^{W}_{+})$ 同构于 $W$ 的正则表示，因此作为 $W$ 模，$\cohog{*}{\sB}$ 也同构于 $W$ 的正则表示。

\begin{remark} \eqref{Spr corr} 中的目标集可以与幂零轨道上不可约 $G$ 等变局部系统的同构类典范地等同。事实上，对幂零轨道 $\calO\subset\sN$ 上不可约 $G$ 等变局部系统 $\calL$，它在 $e\in\calO$ 处的茎给出中心化子 $C_{G}(e)$ 的一个不可约表示 $\rho$，该表示经由 $A_{e}$ 分解。注意当 $G$ 在固定同源类中变动时，$G$ 等变的概念随之改变。若给 $\Irr(W)$ 补上若干较小 Weyl 群的 $\Irr(W')$，可以把上述单射 \eqref{Spr corr} 扩充为双射。这就是 Lusztig \cite{LuIC} 发现的{\em 广义 Springer 对应}。
\end{remark}

Springer 对定理 \ref{th:Spr corr} 的原始证明使用 $\frg(\FF_{q})$ 上的三角和；当 $k$ 的特征为零时，他的证明还用到向有限域的归约。下述 Borho 与 MacPherson 的定理 \cite{BM} 可以对所有大特征或特征零的基域 $k$ 直接证明 Springer 对应。为陈述它，需要可构造 $\Ql$ 层（复形）与反常层的语言，参见标准文献 \cite{BBD} 与本卷中 de Cataldo 的讲义 \cite{dC}。

\begin{theorem}\label{th:BM} 复形 $\calS:=\bR\pi_{*}\Ql[\dim\sN]$ 是 $\sN$ 上的反常层，其自同态环典范同构于群代数 $\Ql[W]$。特别地，$W$ 作用在 $\bR\pi_{*}\Ql$ 的各茎上，即 $W$ 作用在 $\cohog{*}{\sB_{e}}$ 上，对一切 $e\in\sN$。
\end{theorem}

下面概述 $W$ 在 $\bR\pi_{*}\Ql[\dim\sN]$ 上作用的三种构造。

\subsubsection{经由中间扩充的构造}\label{sss:mid} 这一构造（更确切地说，是把 $\frg$ 换成 $G$ 的版本）属于 Lusztig \cite[\S3]{LuGreen}。Springer 纤维的维数公式 \eqref{Spr dim} 蕴含
\begin{itemize}
\item 映射 $\pi:\tN\to\sN$ 是半小的。\footnote{对不可约簇的真满射 $f:X\to Y$，若对任意 $d\geq 1$，集合 $\{y\in Y|\dim f^{-1}(y)\geq d\}$ 在 $Y$ 中的余维至少为 $2d$（相应地 $2d+1$），则称 $f$ 为{\em 半小}（相应地{\em 小}）映射。}
\end{itemize}
对 $\tB_{X}$ 的维数，公式 \eqref{Spr dim} 有适用于一切 $X\in \frg$ 的推广。用这个公式可以证明
\begin{itemize}
\item 映射 $\pi_{\frg}:\tg\to\frg$ 是小映射。
\end{itemize}
反常层理论中有一个众所周知的事实：$\pi_{\frg}$ 的小性（连同 $\tg$ 光滑这一事实）蕴含 $\calS_{\frg}:=\bR\pi_{\frg,*}\Ql[\dim\frg]$ 是反常层，且是它在 $\frg$ 任一稠密开子集上的限制的中间扩充。在正则半单轨迹 $\frg^{\rs}\subset \frg$ 上方，$\pi_{\frg}$ 是 $W$-挠子，因此 $\calS_{\frg}|_{\frg^{rs}}$ 是次数移位 $-\dim\frg$ 的局部系统，并带有 $W$ 的作用。由中间扩充的函子性，$\calS_{\frg}$ 带有 $W$ 的作用。取 $\calS_{\frg}$ 的茎，就得到 $W$ 在 $\cohog{*}{\tB_{X}}$ 上的作用，对一切 $X\in\frg$。

特别地，对幂零元 $e$，我们得到 $W$ 在 $\cohog{*}{\sB_{e}}=\cohog{*}{\tB_{e}}$ 上的作用，因为 $\tB_{e}$ 与 $\sB_{e}$ 有相同的约化结构。这正是 Springer 在其原始论文 \cite{Springer} 中定义的作用，{\em 它与我们的作用相差与 $W$ 的符号特征作张量}。
